\section{Introduction}
\label{sec:introduction}

The deployment of Large Language Models (LLMs) has transitioned from passive conversational assistants to autonomous agentic architectures capable of invoking external APIs, executing database queries, and modifying filesystem structures~\cite{yu2024toolprivbench, zhang2024autocoderover}. In mission-critical environments, adhering to the foundational security principle of \textit{Least Privilege} (PoLP) is paramount: an agent should be granted the minimal set of privileges necessary to accomplish its designated task.

Recent studies on English benchmarks, such as \textsc{ToolPrivBench}~\cite{yu2024toolprivbench} and \textsc{Fortis}~\cite{zhang2024fortis}, demonstrate that even frontier models (e.g., GPT-4o, Claude 3.5) suffer from \textit{over-privileged tool selection}. When presented with multiple tools that satisfy a query, agents default to over-privileged tools in up to 38\% of unconstrained interactions. Crucially, when transient execution failures (e.g., HTTP 504 Gateway Timeouts) are injected into the minimal tool, models exhibit \textit{panic privilege escalation}, escalating to unrestricted administrative tools in over 60\% of cases.

Despite these alarming findings, existing benchmarks and defense mechanisms suffer from two major blindspots:
\begin{enumerate}
    \item \textbf{Linguistic Monoculture:} Existing evaluations operate exclusively on high-resource English text. Southeast Asian languages, representing over 680 million speakers, display complex morphological variations and honorific speech levels (e.g., Javanese \textit{Ngoko} vs. \textit{Krama}) that introduce semantic dispersion. How multilingual foundation models (e.g., SEA-LION~\cite{aisingapore2024sealion}) behave under cross-lingual tool authorization remains entirely unmeasured.
    \item \textbf{Inefficacy of Prompt-Based Defenses:} Existing attempts to curb over-privilege rely on system prompt instructions (e.g., \textit{``Please pick the least-privileged tool''}). As demonstrated by prior work, stochastic prompt-level controls degrade under stress and fail completely when exposed to transient runtime errors.
\end{enumerate}

To address these gaps, this paper makes three key contributions:
\begin{itemize}
    \item \textbf{Cross-Lingual Vulnerability Benchmark (ToolPrivBench-SEA):} We adapt and curate a culturally grounded evaluation set of tool-calling scenarios across English, Indonesian, and Javanese, annotated with formal permission scopes and verified by independent native annotators.
    \item \textbf{Deterministic Least-Privilege Guardrail:} We design a mechanistic verification layer based on mathematical Lattice Theory that intercepts tool emissions, deterministically validates parameter scopes, and performs automatic down-casting to the minimal sufficient equivalent.
    \item \textbf{Empirical Fault Robustness:} We benchmark open-weight multilingual models and prove that our deterministic guardrail prevents unauthorized privilege escalation under simulated transient API faults without degrading task completion utility.
\end{itemize}
